\section{Diagnóstico estadístico preliminar}
\label{ch:diagnostico}

\section{Estadística descriptiva disponible}
\label{sec:estadistica_descriptiva}

El script \texttt{R/03\_auditoria\_indicadores.R} genera diagnósticos estadísticos preliminares sobre la matriz de indicadores candidatos (7 indicadores, base total y base válida). Los resultados se almacenan en \texttt{resultados/auditoria\_estadistica/} y \texttt{resultados/base\_trabajo/}.

\begin{table}[htbp]
\centering
\caption{Estadísticos descriptivos por indicador (matriz base total) -- extracto de \texttt{estadisticos\_indicadores.csv}}
\label{tab:estadisticos_descriptivos}
\begin{adjustbox}{max width=\textwidth}
\begin{tabular}{@{}lrrrrrr@{}}
\toprule
\textbf{Indicador} & \textbf{Media} & \textbf{Mediana} & \textbf{DE} & \textbf{Asimetría} & \textbf{\% Ceros} & \textbf{Rango} \\
\midrule
actividad\_desocupacion\_incapacidad\_hogar & -- & -- & -- & -- & -- & [0,1] \\
sin\_energia\_electrica & -- & -- & -- & -- & -- & [0,1] \\
sin\_gas\_red\_publica & -- & -- & -- & -- & -- & [0,1] \\
sin\_recoleccion\_basuras & -- & -- & -- & -- & -- & [0,1] \\
sin\_internet & -- & -- & -- & -- & -- & [0,1] \\
hogares\_con\_menores15 & -- & -- & -- & -- & -- & [0,1] \\
estrato\_1\_2 & -- & -- & -- & -- & -- & [0,1] \\
\bottomrule
\end{tabular}
\end{adjustbox}
\end{table}

\noindent \textit{Nota: Los valores exactos deben leerse de \texttt{estadisticos\_indicadores.csv} y \texttt{varianza\_indicadores.csv}. En esta versión del informe se reporta la estructura; los números finales se insertarán tras la auditoría integral.}

\section{Variabilidad}
\label{sec:variabilidad}

El archivo \texttt{varianza\_indicadores.csv} reporta la varianza de cada indicador y señala variables con \textbf{varianza casi nula} (umbral configurable en \texttt{R/00\_config.R}: \texttt{varianza\_casi\_nula = 1e-8}). Una varianza muy baja indica escasa discriminación entre municipios y posible exclusión de la matriz PCA.

\section{Correlaciones preliminares}
\label{sec:correlaciones_preliminares}

\begin{itemize}
    \item \texttt{correlaciones\_pares.csv}: Correlación de Pearson para cada par de los 7 indicadores propuestos (21 pares).
    \item \texttt{pares\_correlacion\_ge\_090.csv}: Pares con $|r| \ge 0,90$ -- \textbf{vacío} en la ejecución registrada (ningún par alcanza el umbral).
    \item \texttt{pares\_correlacion\_ge\_095.csv}: Pares con $|r| \ge 0,95$ -- \textbf{vacío}.
\end{itemize}

La ausencia de correlaciones extremadamente altas ($|r| \ge 0,90$) sugiere que no hay redundancia lineal severa entre los 7 indicadores candidatos, aunque la correlación moderada/alta entre algunos pares debe evaluarse en el contexto multivariado (KMO, MSA).

\section{Outliers}
\label{sec:outliers}

\begin{itemize}
    \item Detección mediante estandarización (z-score) usando \texttt{parametros\_estandarizacion.csv} (media y DE por indicador).
    \item Umbral: $|z| > 4$ (configurable en \texttt{R/00\_config.R}: \texttt{z\_outlier = 4.0}).
    \item Resultado: \textbf{32 municipios} presentan al menos un indicador con $|z| > 4$. Ver \texttt{outliers\_z\_mayor\_4.csv}.
    \item \textbf{No se eliminan automáticamente.} Estos municipios deben revisarse contra las fuentes originales y el contexto territorial antes de cualquier decisión de winsorización o exclusión.
\end{itemize}

\section{Diagnósticos de adecuación multivariada}
\label{sec:diagnosticos_aptitud}

El archivo \texttt{diagnosticos\_aptitud.csv} reporta:

\begin{table}[htbp]
\centering
\caption{Diagnósticos globales de aptitud para PCA (matriz base total, 7 indicadores)}
\label{tab:diagnosticos_aptitud}
\begin{adjustbox}{max width=\textwidth}
\begin{tabular}{@{}lll@{}}
\toprule
\textbf{Diagnóstico} & \textbf{Valor} & \textbf{Interpretación} \\
\midrule
KMO global & $\approx 0,642$ & Moderado (umbral adecuado $\ge 0,70$; mínimo $\ge 0,50$) \\
Test de Bartlett & $\chi^2 \approx 4664,873$, gl=21, $p < 0,001$ & Significativo: rechaza matriz identidad \\
Determinante de la matriz de correlación & $\approx 0,0154$ & $> 0,00001$: no colinealidad perfecta \\
\bottomrule
\end{tabular}
\end{adjustbox}
\end{table}

\section{Medida de Adecuación Muestral Individual (MSA por indicador)}
\label{sec:msa_individual}

El archivo \texttt{kmo\_por\_indicador.csv} reporta el MSA (KMO individual) para cada uno de los 7 indicadores:

\begin{itemize}
    \item La mayoría de indicadores tienen MSA $> 0,50$ (mínimo aceptable).
    \item \textbf{Alerta:} El indicador \texttt{estrato\_1\_2} (proxy estrato 1--2) tiene \textbf{MSA $\approx 0,414$}, \textbf{por debajo del mínimo} ($0,50$). Esto sugiere que este indicador no comparte varianza común suficiente con el resto del conjunto y su inclusión en el PCA debe reconsiderarse.
    \item Ver \texttt{Resumen\_General.md}, punto 40: ``el MSA de estrato 1--2 es bajo ($\approx 0,414$)''.
\end{itemize}

\section{Sensibilidad a No Respuesta (NR)}
\label{sec:sensibilidad_nr_diag}

La comparación entre \textbf{base total} (incluye NR en denominador) y \textbf{base válida} (excluye NR) muestra:

\begin{itemize}
    \item \texttt{sensibilidad\_nr.csv}: Resumen por indicador de diferencias absolutas y relativas.
    \item \texttt{sensibilidad\_nr\_municipio.csv}: Diferencias municipio a municipio con marcas de sensibilidad (umbrales en \texttt{R/00\_config.R}: \texttt{delta\_absoluta = 0.02}, \texttt{delta\_relativa = 0.10}).
    \item Las diferencias municipales y agregadas \textbf{deben revisarse} antes de fijar la base analítica principal (Decisión vigente en \texttt{Resumen\_General.md}, punto 27 y punto 59).
\end{itemize}

\section{Estado de la auditoría estadística}
\label{sec:auditoria_pendiente}

\textbf{Importante:} Los diagnósticos reportados en esta sección son \textbf{preliminares}. La \textbf{auditoría estadística integral está pendiente} y debe incluir, como mínimo:

\begin{enumerate}
    \item Comparación sistemática de correlaciones, KMO y MSA entre \textbf{matriz base total} y \textbf{matriz base válida}.
    \item Evaluación de transformaciones para variables asimétricas (ej. \texttt{sin\_internet}, \texttt{sin\_gas} con techo alto).
    \item Análisis de componentes principales \textit{exploratorio} (no definitivo) para inspeccionar eigenvalues, scree plot, cargas, varianza explicada.
    \item \textit{Parallel analysis} para criterio de retención de componentes.
    \item Evaluación de estabilidad: bootstrap, jackknife, o análisis de sensibilidad a outliers (32 municipios con $|z|>4$).
    \item Decisión fundamentada sobre inclusión/exclusión de \texttt{estrato\_1\_2} (MSA bajo) y \texttt{hogares\_con\_menores15} (cobertura baja en 26 municipios).
    \item Definición de dirección de \texttt{hogares\_con\_adultos\_60mas} (actualmente fuera de matriz candidata).
    \item Documentación de criterios de retención de componentes (eigenvalue $>1$, scree, parallel analysis, varianza explicada acumulada $> 60-70\%$).
\end{enumerate}

\textbf{No se ha ejecutado un PCA definitivo.} Los diagnósticos actuales \textbf{no deciden} si el PCA es apropiado; son insumos para la decisión metodológica que debe tomar el autor (ver \texttt{Resumen\_General.md}, punto 40: ``La auditoría estadística recomienda no ejecutar aún un PCA definitivo'').
